\chapter{Conclusioni e Sviluppi Futuri}
\label{ch:conclusioni}

\section{Sintesi dei risultati ottenuti}
\label{sec:sintesi_risultati}

Il presente lavoro di tesi ha progettato, implementato e valutato empiricamente un framework basato su architetture neurali profonde per la compressione, la ricostruzione e la generazione stocastica di matrici di correlazione finanziarie realistiche, applicato a un paniere di $N=362$ titoli dell'indice S\&P 500 osservati lungo un orizzonte ventennale (2000-2020, Capitolo~\ref{ch:preprocessing}).

Il percorso metodologico ha seguito un'evoluzione incrementale, motivata empiricamente passo dopo passo. Il primo tentativo di ricostruzione, condotto direttamente nel dominio della matrice di correlazione grezza, ha evidenziato che le matrici ricostruite non rispettavano in modo affidabile il vincolo di semidefinita positività (PSD): un limite non accettabile per una matrice di correlazione, che ha motivato il passaggio, adottato in via definitiva, alla ricostruzione della decomposizione di Cholesky (Capitoli~\ref{ch:modelli_generativi} e~\ref{ch:architetture}), garantendo la PSD per costruzione matematica indipendentemente dalla qualità dell'apprendimento.

Sul piano della pura ricostruzione storica (Capitolo~\ref{ch:analisi_compressione}), i risultati sono coerenti con il teorema di \citet{bourlard1988autoassociation}: PCA e Autoencoder Lineare convergono, a parità di dimensione latente $K$, verso lo stesso limite di errore, stabilendo il \emph{benchmark} lineare. L'Autoencoder Profondo, grazie alla non-linearità e all'omogeneizzazione statistica ottenuta con il campionamento a mescolamento casuale (Capitolo~\ref{ch:preprocessing}), supera nettamente tale limite a ogni $K$, senza mostrare segni di \emph{overfitting}. Il Variational Autoencoder, nella parametrizzazione Cholesky con loss basata sulla Divergenza di Jeffreys (modello \texttt{VAE\_20dim\_cholesky\_06\_loss}, adottato come riferimento definitivo), si colloca in una posizione intermedia: pur restando meno accurato dell'Autoencoder Profondo deterministico, supera sia la PCA sia l'Autoencoder Lineare a $K=20$ --- un risultato che ridimensiona, rispetto a possibili aspettative iniziali, il costo in termini di fedeltà di ricostruzione imposto dalla regolarizzazione variazionale.

È tuttavia sul fronte generativo che il VAE mostra il proprio valore distintivo (Capitolo~\ref{ch:generazione_stocastica}). Sfruttando la continuità e la densità dello spazio latente indotte dalla Divergenza di Kullback-Leibler, è stato possibile costruire una procedura di generazione --- un \emph{block-bootstrap} storico degli incrementi latenti giornalieri --- capace di produrre, a partire dallo stato di mercato osservato oggi, una distribuzione di scenari sintetici della struttura di correlazione a un orizzonte nominale di 21 giorni: non una previsione puntuale, bensì una stima della possibile variabilità delle correlazioni a quell'orizzonte, valutabile in ottica di rischiosità. Le matrici sintetiche così generate, verificate rispetto ai cinque fatti stilizzati delle matrici di correlazione finanziarie (Capitolo~\ref{ch:stylized_facts}), non mostrano violazioni: la distribuzione dei coefficienti \emph{pairwise} resta spostata verso valori positivi, lo spettro degli autovalori mantiene la stessa ripartizione tra effetto mercato, effetti settoriali e rumore, la proprietà di Perron-Frobenius è preservata con similarità coseno superiore a 0.98 su tutti i modi principali, e sia la struttura gerarchica sia la topologia \emph{scale-free} del Minimum Spanning Tree risultano qualitativamente conservate. Un Autoencoder Profondo puramente deterministico, per quanto superiore nella ricostruzione storica, non è stato impiegato per lo stesso compito: come noto in letteratura \citep{kingma2013vae, doersch2016vae, burgess2018understanding}, in assenza della regolarizzazione variazionale lo spazio latente di un AE classico non è vincolato ad alcuna distribuzione a priori e non possiede le proprietà di continuità e densità necessarie a un campionamento significativo --- un limite che nel presente lavoro non è stato verificato empiricamente, ma assunto come premessa metodologica (Capitolo~\ref{ch:modelli_generativi}, \S\ref{sec:deep_ae}).

\paragraph{Differenze e approfondimenti rispetto al lavoro di riferimento.} L'impostazione generale di questo lavoro riprende dichiaratamente quella di \citet{caprioli2024ptfvae}, indicata fin dal Capitolo~\ref{ch:introduzione}, \S\ref{sec:problema_ricerca}, come il precedente più affine: gli autori addestrano un VAE con spazio latente bidimensionale su un dataset di $206$ matrici di correlazione mensili di $44$ indici azionari (finestre mobili di $100$ mesi, 1997--2022; ripartizione casuale $70\%/30\%$ fra addestramento e validazione), ne confrontano l'errore di ricostruzione con quello di un autoencoder deterministico e di autoencoder lineari, ne verificano l'output rispetto ai cinque fatti stilizzati di \citet{marti2019corrgan} e ricampionano tramite \emph{bootstrap} le variazioni delle coordinate latenti per stimare la distribuzione del \emph{Value at Risk} (VaR) di un portafoglio creditizio. Rispetto a tale impianto, assunto come punto di partenza, il presente lavoro presenta differenze e approfondimenti sotto i profili seguenti, che è utile esplicitare in chiusura.
\begin{itemize}
    \item \textbf{Dataset più ampio in sezione trasversale e più denso nel tempo.} Al posto di $44$ indici aggregati osservati a frequenza mensile, il paniere è costituito da $N=362$ singoli titoli dell'S\&P 500 a frequenza giornaliera lungo un ventennio (Sezione~\ref{sec:dataset_desc}), da cui la finestra mobile di $724$ giorni con passo di $10$ giorni estrae $461$ matrici, ridotte a $412$ dal filtro di Perron-Frobenius (Sezione~\ref{sec:purging}), ferma restando la distorsione da sopravvivenza discussa nella Sezione~\ref{sec:limiti_modello}. Il problema cambia scala: ogni matrice conta $65703$ coefficienti nella parametrizzazione di Cholesky, contro i $1936$ nodi di ingresso ($44 \times 44$, corrispondenti a $946$ coefficienti distinti) del VAE di riferimento. Il passaggio dagli indici ai singoli titoli rende inoltre più stringente, e più significativa, la verifica dei fatti stilizzati legati alla granularità del paniere --- la struttura gerarchica settoriale e la topologia \emph{scale-free} dell'MST --- proprietà che gli stessi autori del lavoro di riferimento osservano su indici aggregati, segnalando, a proposito della concentrazione spettrale e della struttura gerarchica, la differenza rispetto ai singoli titoli analizzati da \citet{marti2019corrgan}.
    \item \textbf{\emph{Sanity check} della pipeline rispetto al \emph{benchmark} lineare e riscontro empirico del teorema di Bourlard e Kamp.} Il confronto con i modelli lineari, che nel lavoro di riferimento misura l'errore di ricostruzione con il solo MSE a due e tre dimensioni latenti e richiama dalla letteratura l'equivalenza fra autoencoder lineare e PCA senza verificarla empiricamente, è qui reso sistematico sui tre modelli deterministici e impiegato come verifica di coerenza della catena di addestramento e valutazione rispetto al \emph{benchmark} lineare analitico: PCA e Autoencoder Lineare sono implementati e addestrati separatamente sulla medesima parametrizzazione (Sezioni~\ref{sec:setup_pca} e~\ref{sec:setup_linear_ae}), la loro convergenza verso lo stesso errore a parità di $K$ --- il contenuto del teorema di \citet{bourlard1988autoassociation} --- è misurata numericamente su una griglia di quattro dimensioni latenti $K \in \{3, 10, 20, 100\}$ (Tabella~\ref{tab:err_test_cholesky}), e la valutazione avviene su un \emph{Test Set} tenuto fuori dall'addestramento e distinto dal \emph{Validation Set} usato per il \emph{checkpointing} (Sezione~\ref{sec:purging}). La stessa griglia, estesa all'Autoencoder Profondo, ha permesso di caratterizzare il vantaggio della non-linearità come vantaggio di efficienza rappresentativa --- a $K=10$ l'AE eguaglia la qualità che la PCA raggiunge solo a $K=100$ --- e di collocare a $K=20$ il gomito della curva di ricostruzione dell'AE (Sezione~\ref{subsec:metriche_errore}).
    \item \textbf{Architettura, funzione di perdita e protocollo di addestramento.} A fronte dell'unica architettura del VAE di riferimento, a tre strati per ramo ($1936 \to 512 \to 250 \to 2$ e speculare) e addestrata per $80$ epoche con \emph{loss} $\mathrm{MSE} + \beta\,\mathrm{KL}$, le reti profonde di questo lavoro --- sperimentate su più dimensioni latenti e, in una prima fase, su due parametrizzazioni dell'input (Sezione~\ref{sec:flattening}) --- adottano un \emph{design} a imbuto $65703 \to 2048 \to 1024 \to 512 \to K$ con LeakyReLU, addestrato per $1000$ epoche (valore fissato per il VAE e tipico per l'Autoencoder Profondo) con Adam, \emph{scheduler} a \emph{Cosine Annealing} e \emph{Model Checkpointing} sul minimo della \emph{Validation Loss} (Sezione~\ref{sec:setup_deep_ae}, Tabella~\ref{tab:vae_iperparametri}); il VAE adotta un obiettivo di tipo VAE con peso unitario ($\beta = 1$) e sostituisce l'MSE con una \emph{Reconstruction Loss} basata sulla Divergenza di Jeffreys, valutata nel dominio delle matrici di correlazione attraverso la ricomposizione di Cholesky mantenuta all'interno del grafo differenziabile (Sezione~\ref{sec:setup_vae}), e la stabilità del suo addestramento è documentata dalla decomposizione della \emph{loss} nei due termini di ricostruzione e di regolarizzazione (Sezione~\ref{subsec:addestramento_vae}).
    \item \textbf{Vincolo di semidefinita positività incorporato nella rappresentazione.} Il VAE di riferimento ricostruisce direttamente i coefficienti della matrice di correlazione (ingresso e uscita a $44 \times 44$ nodi), parametrizzazione che, come discusso nella Sezione~\ref{sec:flattening}, non offre di per sé alcuna garanzia strutturale di semidefinita positività. Qui l'input di tutti i modelli, ricostruttivi e generativi, è il \emph{flattening} del fattore triangolare di Cholesky, e la matrice viene ricomposta come $\hat{\mathbf{L}}\hat{\mathbf{L}}^T$ seguita da rinormalizzazione: la semidefinita positività è così garantita per costruzione, per la PCA come per le reti, e viene confermata a posteriori su tutti i $1000$ scenari generati (autovalore minimo strettamente positivo, Sezione~\ref{subsec:verifica_aggregata}).
    \item \textbf{Batteria di metriche di errore puntuali e topologiche.} La bontà della ricostruzione non è misurata dal solo MSE, come nel lavoro di riferimento, ma da tre metriche di errore spaziali (MSE, MAE, norma di Frobenius) e da cinque metriche topologiche di similarità fra il Minimum Spanning Tree della matrice originale e quello della matrice ricostruita (\emph{edge overlap}, indice di Jaccard, sovrapposizione dei nodi a maggiore \emph{betweenness}, scarto della lunghezza media dei cammini pesata e non pesata; Sezione~\ref{sec:metriche_convergenza}). L'ordinamento fra modelli si è rivelato invariante rispetto alla metrica scelta, e le metriche topologiche hanno messo in luce un fenomeno che gli errori puntuali lasciano soltanto intuire: la severità del crollo della struttura di rete ricostruita dai modelli lineari alle basse dimensioni latenti (Sezioni~\ref{subsec:metriche_errore} e~\ref{subsec:metriche_mst}).
    \item \textbf{Determinazione empirica della dimensionalità effettiva dello spazio latente.} Laddove il lavoro di riferimento fissa \emph{a priori} uno spazio latente bidimensionale, privilegiandone la leggibilità, qui la dimensione nominale è stata scelta sul gomito della curva di ricostruzione ($K=20$) e la dimensionalità \emph{effettiva} è stata poi misurata sullo spettro di attività delle dimensioni latenti: un nucleo di sei fattori fortemente attivi (varianza latente $\geq 0.5$), invariante al variare di $K \in \{10, 20, 30\}$, entro un sottospazio attivo di circa sei--nove dimensioni, con collasso a posteriori delle dimensioni in eccesso (Sezioni~\ref{subsec:vae_vs_k} e~\ref{subsec:struttura_latente}, Tabella~\ref{tab:vae_vs_k_test}). La dimensionalità della \emph{manifold} delle correlazioni emerge così --- nei limiti dell'intervallo di $K$ esplorato --- come una proprietà dei dati che il modello individua da sé, più che come un iperparametro imposto.
    \item \textbf{Arricchimento del dataset con la classificazione settoriale GICS.} Il paniere è stato corredato della classificazione GICS di ciascun titolo, ricostruita --- in assenza di un \emph{provider} dati strutturato --- tramite un modello linguistico di grandi dimensioni interrogato con un prompt \emph{few-shot} ($360$ titoli su $362$ classificati; Sezione~\ref{sec:dataset_desc}, Tabella~\ref{tab:gics_distribution}), con il limite di validazione dichiarato nella Sezione~\ref{sec:limiti_modello}. Tale informazione rende leggibile la struttura gerarchica del mercato (dendrogrammi con foglie colorate per settore, Sezione~\ref{sec:struttura_gerarchica}) e, soprattutto, ne consente la misura quantitativa attraverso l'indicatore introdotto al punto seguente; in \citet{caprioli2024ptfvae}, i cui asset sono essi stessi indici di mercato e settoriali, la struttura gerarchica è osservata su \emph{heatmap} e dendrogrammi di singole matrici scelte casualmente, senza che l'informazione settoriale sia tradotta in un indicatore.
    \item \textbf{Validazione quantitativa ed estesa dei fatti stilizzati, con un indicatore definito \emph{ad hoc} per la struttura gerarchica.} La verifica dei cinque fatti stilizzati, che nel lavoro di riferimento è condotta per via prevalentemente grafica su matrici decodificate da una griglia di punti dello spazio latente, è qui tradotta in un protocollo quantitativo: per ciascuno dei $1000$ scenari generati si calcola un indicatore scalare per ogni proprietà --- media dei coefficienti \emph{pairwise}, autovalore di mercato e numero di modi oltre il limite $\lambda_+$ di Marchenko-Pastur, componente minima del \emph{market mode}, rapporto intra/inter-settore $R$ (Equazione~\ref{eq:intra_inter}) ed esponente $\gamma$ della legge di potenza dell'MST --- con un criterio di conformità esplicito e con il confronto sia con la matrice odierna sia con l'intervallo osservato sulle $412$ matrici reali (Tabella~\ref{tab:sf_aggregate}). Il rapporto $R$, in particolare, è stato definito appositamente per rendere misurabile ciò che il dendrogramma mostra solo visivamente. A monte, il significato dei cinque fatti è stato saggiato su un \emph{benchmark} di rumore gaussiano puro di pari dimensioni (Sezione~\ref{sec:benchmark_rmt}), e sull'intero bacino di matrici reali il filtro di Perron-Frobenius ha isolato le $49$ finestre iniziali che violano la proprietà, dandone un'interpretazione economica (Sezione~\ref{sec:purging}); a valle, l'ispezione di un singolo scenario aggiunge la similarità coseno dei primi autovettori e il confronto diretto di spettro, dendrogramma e distribuzione di grado (Sezione~\ref{subsec:verifica_dettaglio}).
    \item \textbf{\emph{Block-bootstrap} latente a passo giornaliero e validazione diretta degli scenari generati.} Il \emph{block-bootstrap} delle coordinate latenti, che il lavoro di riferimento applica --- insieme a un \emph{bootstrap} semplice --- alle differenze mensili delle due coordinate latenti su un orizzonte annuale, è stato adattato a una traiettoria latente giornaliera ($4124$ osservazioni, codificate a passo unitario) in uno spazio a $20$ dimensioni, con blocchi di $21$ incrementi consecutivi e orizzonte nominale di un mese (Sezione~\ref{sec:block_bootstrap}). I $1000$ scenari così ottenuti sono gli stessi sottoposti alla verifica aggregata, sicché la validazione dei fatti stilizzati riguarda direttamente la distribuzione di scenari destinata all'uso in ottica di rischiosità, laddove nel lavoro di riferimento i fatti stilizzati sono verificati su punti dello spazio latente distinti dal campione \emph{bootstrap} impiegato per la stima del rischio.
\end{itemize}
Restano invece fuori dal perimetro di questo lavoro due elementi che in quello di riferimento costituiscono il punto di arrivo: la traduzione della distribuzione degli scenari in una distribuzione di misure di rischio di portafoglio --- indicata come primo sviluppo futuro nella Sezione~\ref{sec:sviluppi_futuri} --- e l'interpretazione economico-finanziaria delle singole coordinate latenti (ad esempio in termini di autovalore di mercato), che uno spazio bidimensionale rende più diretta e che uno spazio a sei--nove fattori attivi rende necessariamente più articolata.

\section{Limiti del modello attuale}
\label{sec:limiti_modello}

Diversi limiti metodologici, emersi nel corso del lavoro, meritano di essere dichiarati esplicitamente; il primo di essi, relativo al dataset di partenza, merita un rilievo particolare, perché condiziona l'input di tutte le analisi svolte.

In primo luogo, il dataset di partenza è soggetto a un \emph{selection bias} nella forma del \emph{survivorship bias}. Il requisito di quotazione ininterrotta lungo l'intero ventennio (Capitolo~\ref{ch:preprocessing}), indispensabile per disporre di un panel bilanciato e privo di dati mancanti, esclude per costruzione le aziende fallite e/o uscite dall'indice nel periodo considerato. Le 362 serie storiche analizzate descrivono quindi il comportamento dei soli titoli sopravvissuti, con il rischio di sottorappresentare le dinamiche di correlazione tipiche delle fasi di dissesto societario e di ricomposizione dell'indice; poiché tali episodi si concentrano proprio nei periodi di stress dei mercati, le matrici empiriche impiegate per l'addestramento --- e di conseguenza gli scenari sintetici da esse derivati --- possono risultare conservativi nella descrizione delle situazioni di crisi. Va sottolineato che questo è il limite più forte dell'intero lavoro, perché agisce a monte di ogni fase successiva: le matrici empiriche estratte a finestra mobile (Capitolo~\ref{ch:preprocessing}), i modelli addestrati e confrontati (Capitoli~\ref{ch:architetture} e~\ref{ch:analisi_compressione}), la verifica dei fatti stilizzati e gli scenari sintetici generati (Capitolo~\ref{ch:generazione_stocastica}) discendono tutti dallo stesso panel di soli titoli sopravvissuti, e nessuna elaborazione a valle --- per quanto accurata --- può reintrodurre un'informazione assente dall'input. Ne consegue che ogni risultato presentato, dalle metriche di ricostruzione alla conformità degli scenari ai fatti stilizzati, va letto come riferito al mercato dei titoli sopravvissuti e non all'universo completo dell'indice. Tale distorsione è stata accettata consapevolmente come limite del dataset di partenza, in coerenza con quanto dichiarato nel Capitolo~\ref{ch:introduzione}, e la sua rimozione costituisce una delle direzioni di sviluppo indicate nella Sezione~\ref{sec:sviluppi_futuri}.

In secondo luogo, la validazione dei cinque fatti stilizzati sulle matrici generate (Capitolo~\ref{ch:generazione_stocastica}) riduce ciascuna proprietà a un indicatore scalare quando viene condotta sull'intero insieme dei 1000 scenari (Sezione~\ref{subsec:verifica_aggregata}), mentre l'ispezione visiva completa --- forma dello spettro, sagoma del dendrogramma, topologia dell'MST --- resta limitata a un singolo scenario rappresentativo. A questo si aggiunge il fatto che i 1000 scenari derivano da blocchi di 21 giorni ampiamente sovrapposti, con una numerosità campionaria effettiva inferiore a 1000, e che l'intera verifica è ancorata a un unico stato di mercato di partenza: la conformità documentata riguarda l'orizzonte nominale di 21 giorni dalla fine del campione e non è automaticamente estendibile a stati iniziali strutturalmente diversi.

In terzo luogo, la strategia di campionamento a mescolamento casuale (Capitolo~\ref{ch:preprocessing}), pur affrontando efficacemente il problema del \emph{distribution shift} di regime macroeconomico tra \emph{Training}, \emph{Validation} e \emph{Test Set}, non elimina la sovrapposizione informativa tra finestre di correlazione adiacenti (fino al 98.6\% di osservazioni condivise tra finestre a indice quasi consecutivo). L'assenza di un vincolo di distanza minima tra gli indici assegnati a split diversi rappresenta un potenziale rischio di ottimismo nelle metriche di generalizzazione riportate.

In quarto luogo, il confronto tra modelli non copre in modo sistematico e uniforme tutte le dimensioni latenti $K$ per il VAE (non testato a $K=3$ o $K=100$), ma l'errore di ricostruzione del VAE non decresce monotonicamente con $K$ come invece accade per gli altri tre modelli --- un comportamento coerente con il collasso a posteriori delle dimensioni in eccesso documentato nella Sezione~\ref{subsec:vae_vs_k}.

In quinto luogo, il bacino di addestramento è deliberatamente circoscritto al regime di mercato ``standard'' o cooperativo: l'esclusione delle 49 finestre che non soddisfano la proprietà di Perron-Frobenius (Capitolo~\ref{ch:preprocessing}, Sezione~\ref{sec:purging}) delimita il perimetro di validità del modello, i cui scenari non rappresentano i regimi eccezionali in cui un bene rifugio si disaccoppia dal mercato azionario --- regimi di particolare interesse in ottica di rischiosità, la cui modellazione richiederebbe un campione di finestre anomale più ricco di quello osservato.

Infine, la classificazione settoriale GICS impiegata per le analisi di struttura gerarchica (Capitolo~\ref{ch:stylized_facts}) è stata ottenuta tramite classificazione automatica affidata a un modello linguistico di grandi dimensioni, in assenza di un provider dati strutturato, e non è stata sottoposta a validazione incrociata con fonti ufficiali.

\section{Sviluppi futuri}
\label{sec:sviluppi_futuri}

Il naturale proseguimento di questo lavoro è l'applicazione pratica delle matrici sintetiche generate e verificate alla gestione del rischio di portafoglio. La procedura di \emph{block-bootstrap} latente qui sviluppata consente di generare scenari realistici per poter ottenere una distribuzione di scenari possibili da valutare in ottica di rischiosità: il passo successivo è impiegare tale distribuzione su portafogli reali, confrontando il quadro di rischio che se ne ricava con quello ottenibile dai soli dati storici tradizionali.

Un secondo filone di sviluppo riguarda l'ulteriore rafforzamento della validazione statistica delle matrici generate. La verifica aggregata di conformità sull'intero insieme dei 1000 scenari, presentata nella Sezione~\ref{subsec:verifica_aggregata}, potrebbe essere ripetuta a partire da una pluralità di stati di mercato iniziali distribuiti lungo il ventennio --- e non dal solo stato finale del campione --- così da verificare la stabilità della conformità al variare del regime di partenza, ed estesa dagli indicatori scalari a test formali di aderenza distributiva tra scenari sintetici e matrici reali.

Un terzo filone riguarda il confronto diretto con paradigmi generativi alternativi basati su reti avversarie (GAN), quale CorrGAN \citep{marti2019corrgan}, per quantificare in modo sistematico i vantaggi e gli svantaggi del VAE (stabilità di addestramento, interpretabilità dello spazio latente) rispetto all'approccio avversariale in questo specifico dominio applicativo. Più in generale, l'integrazione di modelli di diffusione, oggi allo stato dell'arte per la generazione condizionata in molti domini, rappresenta un'ulteriore direzione di ricerca promettente.

Un quarto filone consiste nella rimozione del \emph{survivorship bias} discusso nella Sezione~\ref{sec:limiti_modello}, mediante la ricostruzione di un panel a composizione variabile che includa anche le aziende fallite e/o uscite dall'indice nel corso del periodo di osservazione. Un'estensione di questo tipo richiederebbe di superare il vincolo di dimensionalità fissa dell'input --- ad esempio aggiornando l'universo di riferimento a ogni data di stima e calcolando le correlazioni sulle sole coppie di titoli contemporaneamente quotati, con conseguente adattamento della parametrizzazione di Cholesky e dell'architettura delle reti --- ma permetterebbe di quantificare in che misura la struttura di correlazione appresa, e con essa il profilo di rischio degli scenari generati, si modifichi una volta reintegrati gli episodi di dissesto societario.

Infine, l'estensione della pipeline ad altre asset class (valute, materie prime) permetterebbe di verificare la generalità del framework proposto al di là del mercato azionario statunitense a più alta capitalizzazione qui analizzato, così come l'introduzione di una procedura di \emph{purging} esplicito tra i sottoinsiemi di addestramento, validazione e test irrobustirebbe ulteriormente la stima delle capacità di generalizzazione dei modelli.
